\section{Para praticar no IDE: blog da JetBrains e JetBrains Guide}
\label{ap:jetbrains}

Não é gente que fez o Java, mas é onde muita gente aprende a usá-lo no dia
a dia: o blog do IntelliJ IDEA e os tutoriais gratuitos do JetBrains Guide.
São textos práticos, sempre com o IDE aberto, escritos por developer
advocates como Trisha Gee, Mala Gupta, Marit van Dijk, Helen Scott e Igor
Kulakov. Os mais úteis para cada capítulo também aparecem nas seções
\enquote{Onde aprofundar}.

% Todos os itens de jetbrains.bib entram nas listas por tema abaixo.
\nocite{artigo:jetbrains-writing-tests-junit5,
  artigo:jetbrains-migrating-junit4-junit5,
  artigo:jetbrains-testcontainers-spring-boot,
  artigo:jetbrains-ai-agent-flaky-tests,
  curso:jetbrains-guide-writing-junit5-tests,
  curso:jetbrains-guide-writing-tests-spock,
  curso:jetbrains-guide-testing-debugging,
  artigo:jetbrains-pattern-matching-5-examples,
  artigo:jetbrains-evolution-switch-java,
  artigo:jetbrains-project-loom-intellij,
  artigo:jetbrains-java-25-lts,
  artigo:jetbrains-do-you-really-know-java,
  artigo:jetbrains-java-still-relevant,
  artigo:jetbrains-thread-dumps-virtual-threads,
  artigo:jetbrains-debugging-java-intellij,
  artigo:jetbrains-sources-bytecode-debugging,
  artigo:jetbrains-println-debugging,
  artigo:jetbrains-debug-without-breakpoints,
  artigo:jetbrains-intellij-profiler,
  artigo:jetbrains-profiling-parrot-party,
  artigo:jetbrains-get-started-profiling,
  artigo:jetbrains-5-profiling-features,
  curso:jetbrains-guide-javax-jakarta,
  artigo:jetbrains-java-best-practices,
  artigo:jetbrains-code-smells-null,
  artigo:jetbrains-code-smells-mutation,
  artigo:jetbrains-code-smells-if,
  artigo:jetbrains-basic-refactorings,
  artigo:jetbrains-3-ways-refactor,
  curso:jetbrains-guide-refactoring}

\begin{description}
\item[Como foram conferidos] O blog da JetBrains bloqueia acesso
  automatizado; título, autor e data foram lidos no snapshot do Internet
  Archive, e o link impresso é o do blog. O JetBrains Guide foi aberto
  direto.
\item[O que ficou de fora] O post sobre String Templates: o recurso saiu
  em preview no JDK~21 e~22 e foi retirado (o JEP~465, terceira preview,
  está \emph{Closed/Withdrawn}); ensinar a sintaxe hoje seria ensinar algo
  que não existe no Java. Também saíram as pesquisas \emph{State of
  Developer Ecosystem}, com pouco conteúdo de Java, notas de versão do IDE e
  resumos mensais.
\end{description}

\subsection{Testes}
JUnit~5, migração do JUnit~4, Testcontainers e Spock; complementam a
Parte~\ref{parte:testes}.
\ondeachar*{jb-testes}

\subsection{Java moderno}
Pattern matching, \texttt{switch}, Project Loom e o Java~25; complementam
o capítulo de Java moderno.
\ondeachar*{jb-moderno}

\subsection{JVM, depuração e profiling}
Thread dumps com virtual threads, bytecode no depurador e o IntelliJ
Profiler; complementam os capítulos de JVM e de bytecode.
\ondeachar*{jb-jvm}

\subsection{Código limpo, refatoração e Jakarta}
A série \emph{Code Smells}, refatoração com o IDE e a migração de
\texttt{javax} para \texttt{jakarta}, que o Quarkus usa.
\ondeachar*{jb-limpo}
